% streaming_architecture.tex — Video Streaming and Ground Station Architecture
% Intended as a detailed design-rationale section (cross-references sec:groundstation)

\section{Video Streaming and Ground Station Architecture}\label{sec:streaming-arch}

The choice of video transport, concurrency model, and ground station topology has a
disproportionate effect on operator situational awareness in small-UAS SAR operations.
This section analyses the trade-offs that led to the current MJPEG-over-HTTP
architecture and documents the alternatives that were evaluated and rejected.

% ═══════════════════════════════════════════════════════════════════════
\subsection{Streaming Protocol Selection}\label{sec:stream-protocol}

Four candidate protocols were evaluated against five criteria relevant to a
Raspberry~Pi companion computer operating over an ad-hoc WiFi link: glass-to-glass
latency, browser compatibility, dependency footprint on the Pi, firewall
transparency, and implementation complexity.

\begin{table}[ht]
\centering
\caption{Comparison of video streaming protocols for companion-computer UAV telemetry.}
\label{tab:stream-protocols}
\small
\begin{tabular}{@{}lccccl@{}}
\toprule
\textbf{Protocol} & \textbf{Latency} & \textbf{Browser} & \textbf{Pi deps} & \textbf{Firewall} & \textbf{Complexity} \\
\midrule
MJPEG/HTTP        & 50--100\,ms & Native \texttt{<img>} & None (stdlib) & TCP 8090 only & Trivial \\
HLS (H.264)       & 2--6\,s     & Native \texttt{<video>} & FFmpeg required & TCP 80 only & Moderate \\
RTP/RTSP          & 30--80\,ms  & Requires VLC/app & GStreamer or FFmpeg & UDP punch-through & Moderate \\
WebRTC            & 30--60\,ms  & Native \texttt{RTCPeerConnection} & libnice, STUN/TURN & ICE negotiation & High \\
\bottomrule
\end{tabular}
\end{table}

\textbf{MJPEG over HTTP} was selected for the following reasons:

\begin{enumerate}
  \item \textbf{Zero external dependencies.} The server uses only Python's standard-library
        \texttt{http.server} and \texttt{socketserver} modules. No FFmpeg binary, no
        GStreamer pipeline, no native codec library needs to be cross-compiled for
        \texttt{aarch64}. On a minimal Raspberry~Pi~OS image where every additional
        package increases the attack surface and deployment friction, this is a decisive
        advantage.

  \item \textbf{Universal browser decoding.} Every modern browser decodes a
        \texttt{multipart/x-mixed-replace} stream natively when assigned to an
        \texttt{<img src="/stream">} element~\cite{munoz2016mjpeg}. No JavaScript media-source
        extension, no WebSocket shim, and no codec negotiation is required. The same URL
        works on a laptop, a tablet, and a phone without installing any application.

  \item \textbf{SSH tunnel transparency.} During bench testing and early field trials, the
        operator frequently connects to the Pi over an SSH tunnel
        (\texttt{ssh -L 8090:localhost:8090 pi@\ldots}). Because MJPEG is plain HTTP,
        the tunnel carries video without modification. HLS would also tunnel, but
        RTP/RTSP requires UDP forwarding (non-trivial over SSH), and WebRTC's ICE
        negotiation fails entirely inside a tunnel.

  \item \textbf{Acceptable latency.} The MJPEG stream adds approximately one frame of
        encoding delay (the \texttt{cv2.imencode} call takes $<$\,2\,ms for a
        640$\times$480 JPEG at quality~85) plus one TCP segment of network delay.
        Measured glass-to-glass latency on local WiFi is 50--100\,ms, well within the
        requirements for operator-in-the-loop confirmation at the \textsc{verify} state,
        where the drone is hovering and the operator has several seconds to respond.
\end{enumerate}

\textbf{HLS (H.264 over HTTP Live Streaming)} was rejected primarily on latency grounds.
HLS segments video into 2--6\,second chunks; even with low-latency HLS extensions, the
minimum achievable latency is approximately 1\,s~\cite{apple_hls}. For an operator
deciding whether a detected blob is a casualty, a one-second delay is tolerable in
isolation but compounds with inference latency and command round-trip time. HLS also
requires FFmpeg to transcode the raw camera frames into H.264 segments, adding a
dependency that must be compiled with hardware-codec support on the Pi---a fragile
build step that would need to be repeated on every OS update.

\textbf{RTP/RTSP} offers the lowest theoretical latency but requires a dedicated client
application (VLC, GStreamer, or a custom player). This eliminates the possibility of
opening the stream on an arbitrary device---precisely the scenario where a field
coordinator needs to glance at the feed on their phone. Furthermore, RTP uses UDP,
which is blocked by many institutional firewalls and cannot traverse a standard SSH
tunnel without additional tooling (\texttt{socat} or a TURN relay).

\textbf{WebRTC} provides browser-native low-latency video with adaptive bitrate, but its
implementation complexity is disproportionate to the project's needs. A minimal WebRTC
stack requires STUN/TURN server configuration, SDP offer/answer exchange, ICE candidate
gathering, and a signalling channel---typically a WebSocket server. The
\texttt{aiortc} Python library provides a pure-Python implementation, but it pulls in
over 20~transitive dependencies and introduces an asyncio event loop that conflicts with
the synchronous MAVLink polling model used throughout the codebase. For a single-viewer
stream on a local network, WebRTC's peer-to-peer negotiation machinery solves a problem
that does not exist.

% ═══════════════════════════════════════════════════════════════════════
\subsection{Threading Model and Frame Buffering}\label{sec:threading-model}

The ground station server uses \texttt{socketserver.ThreadingMixIn} to spawn a new
thread for each incoming HTTP connection. This design was chosen over three
alternatives:

\begin{itemize}
  \item \textbf{Single-threaded sequential server}: would block the \texttt{/state} JSON
        endpoint while a long-lived \texttt{/stream} connection is active. Since the
        browser polls \texttt{/state} every 400\,ms for telemetry updates, blocking would
        cause the dashboard to appear frozen.
  \item \textbf{asyncio (\texttt{aiohttp})}: Python's \texttt{asyncio} model requires
        all blocking calls to be wrapped in executors. The MAVLink
        \texttt{recv\_match()} call, OpenCV's \texttt{imencode()}, and TFLite's
        \texttt{invoke()} are all synchronous C-extension calls that cannot yield to an
        event loop. Wrapping each in \texttt{loop.run\_in\_executor()} would add
        complexity without measurable benefit at the concurrency level of one or two
        simultaneous browser clients.
  \item \textbf{Multi-process}: forking is unnecessary for I/O-bound HTTP serving and
        would complicate shared state (telemetry, frame buffer) by requiring
        inter-process communication.
\end{itemize}

The critical shared resource is the latest encoded JPEG frame. A \emph{latest-frame-wins}
buffer replaces any unread frame with the newest one, ensuring that the stream always
shows the most recent camera image rather than queuing stale frames behind a
backlog~\cite{gamma1994design}. The implementation is a single \texttt{threading.Lock}
protecting a byte-string reference:

\begin{verbatim}
# Producer (main loop, ~50 Hz):
_, jpeg = cv2.imencode('.jpg', frame,
                       [cv2.IMWRITE_JPEG_QUALITY, 85])
with self._frame_lock:
    self._stream_jpeg = jpeg.tobytes()

# Consumer (HTTP handler, per-client thread):
with self._frame_lock:
    jpeg = self._stream_jpeg
\end{verbatim}

No queue is used. If the consumer reads slower than the producer, intermediate frames are
silently dropped---this is the correct behaviour for live video, where showing the
\emph{latest} frame is always preferable to playing back a delayed queue. The HTTP handler
sleeps for 50\,ms between sends, capping the stream at approximately 20\,fps regardless
of how fast the producer runs. This rate exceeds the inference throughput
(${\sim}$\,5\,fps on the Pi) and is therefore never the bottleneck.

% ═══════════════════════════════════════════════════════════════════════
\subsection{Headless Operation}\label{sec:headless-design}

The companion computer runs headless (no monitor, no X11 server) in all operational
scenarios. OpenCV's default behaviour is to call \texttt{cv2.imshow()}, which crashes
immediately on a system without a display server. Rather than conditionally guarding every
GUI call throughout the codebase, the system applies a monkey-patch at import time when
the \texttt{DISPLAY} environment variable is absent:

\begin{verbatim}
if not os.environ.get('DISPLAY'):
    cv2.imshow    = lambda *a, **k: None
    cv2.namedWindow = lambda *a, **k: None
    cv2.waitKey   = lambda *a, **k: -1
\end{verbatim}

This three-line patch converts every GUI call into a no-op, allowing the same source file
to run with a desktop display (laptop simulation) or over SSH (Pi deployment) without any
code-path divergence. All visual output is redirected to the browser via the MJPEG stream,
making the web dashboard the sole operator interface in production.

Operator input in headless mode is accepted through three parallel channels that converge
on a single command dispatcher: (i)~terminal keyboard polling via a dedicated
\texttt{threading.Thread} that reads single characters from \texttt{stdin}; (ii)~browser
buttons that issue \texttt{POST /command} requests; and (iii)~the same REST endpoint
accessed programmatically (e.g.\ \texttt{curl -X POST \ldots}). All three routes call the
same \texttt{handle\_command()} method, guaranteeing consistent behaviour regardless of
input source~\cite{goodrich2007human}.

% ═══════════════════════════════════════════════════════════════════════
\subsection{Port Allocation and Service Separation}\label{sec:port-alloc}

Three HTTP services coexist on the Pi during a typical flight:

\begin{itemize}
  \item \textbf{Port 8090} --- primary ground station (\texttt{pi\_flight.py}) or passive
        observer (\texttt{passive\_watch.py}). Only one of these runs at a time; they
        share the port because they serve the same role (camera stream plus telemetry) and
        binding the same port prevents accidental co-execution.
  \item \textbf{Port 8091} --- training data capture (\texttt{capture\_training.py}). This
        service runs \emph{alongside} the primary ground station during manual flights to
        record photos and video for model retraining. A separate port avoids contention on
        the camera resource and allows the operator to open both dashboards in adjacent
        browser tabs.
  \item \textbf{Port 5762} --- MAVProxy TCP bridge to Mission Planner. This is not an HTTP
        service but a raw MAVLink TCP stream; it is listed here because it shares the Pi's
        network interface and must not collide with the HTTP ports.
\end{itemize}

Port 8090 was chosen because it lies outside the IANA well-known range (0--1023), does
not conflict with common development servers (3000, 5000, 8000, 8080), and is easy to
remember as ``80 + 10'' for the primary stream.

% ═══════════════════════════════════════════════════════════════════════
\subsection{Bandwidth Estimation}\label{sec:bandwidth}

MJPEG transmits each frame as an independent JPEG image, forgoing inter-frame compression.
The per-frame size depends on resolution, quality factor, and scene complexity. Empirical
measurements on representative outdoor aerial imagery give:

\begin{itemize}
  \item \textbf{640$\times$480, quality 70}: 25--40\,KB per frame $\Rightarrow$ at 10\,fps,
        approximately 2.0--3.2\,Mbps.
  \item \textbf{640$\times$480, quality 85} (current setting): 40--60\,KB per frame
        $\Rightarrow$ at 10\,fps, approximately 3.2--4.8\,Mbps.
  \item \textbf{1456$\times$1088, quality 85}: 120--180\,KB per frame $\Rightarrow$ at
        5\,fps, approximately 4.8--7.2\,Mbps.
\end{itemize}

On a dedicated 802.11n WiFi link (theoretical 72\,Mbps at MCS\,7, practical
${\sim}$30\,Mbps), the current configuration consumes less than 15\% of available
bandwidth. For deployment over a cellular (4G) backhaul---where uplink bandwidth may be
limited to 5--10\,Mbps---the quality factor would need to be reduced to 50--60 or the
resolution halved, both of which are single-parameter changes in \texttt{cv2.imencode}.

The bandwidth penalty relative to H.264 is significant: an equivalent H.264 stream at the
same resolution and frame rate would require approximately 1--2\,Mbps due to inter-frame
prediction and entropy coding~\cite{wiegand2003h264}. This 3--5$\times$ overhead is
the principal cost of the MJPEG approach, accepted in exchange for the dependency and
latency advantages described in Section~\ref{sec:stream-protocol}. For the project's
operational context---a short-range WiFi link within a few hundred metres---the overhead is
well within the available capacity.

% ═══════════════════════════════════════════════════════════════════════
\subsection{Ground Station Alternatives Considered}\label{sec:gs-alternatives}

Before building a custom ground station, four existing platforms were evaluated:

\begin{enumerate}
  \item \textbf{QGroundControl}~\cite{qgroundcontrol2020}: a mature, cross-platform GCS
        with MAVLink telemetry, waypoint planning, and video streaming via RTSP. However,
        it cannot display custom AI detection overlays, does not support GPS-projected
        target clustering, and its plugin architecture requires C++ development. Integrating
        the project's Python-based perception pipeline would have required either an IPC
        bridge (adding latency and a failure mode) or reimplementing the detection logic
        in C++.

  \item \textbf{MAVProxy with map module}: a lightweight Python GCS that provides a
        console and a 2D map. It has no video display capability whatsoever, making it
        unsuitable for an operator who needs to visually confirm detections before
        authorising descent.

  \item \textbf{Custom desktop application (PyQt/Tkinter)}: would provide full control
        over the UI but requires a display server, which is unavailable on the Pi. Running
        the desktop app on the laptop and streaming video to it recreates the same
        streaming problem without solving it. Furthermore, desktop frameworks are not
        accessible from a phone or tablet in the field.

  \item \textbf{DroneKit-Android}: an unmaintained SDK (last commit 2019) for building
        Android GCS applications. Its abandonment, combined with the Android-only
        restriction, made it unsuitable for a project that needed to run on any device
        with a browser.
\end{enumerate}

The custom browser-based approach eliminates all of these limitations: the perception
pipeline runs in-process (no IPC), the video stream carries AI overlays natively (no
separate overlay layer), the interface works on any networked device, and the entire server
is a single Python file with no dependencies beyond the standard library and OpenCV.
